\section*{Introduction}

Information is the commodity that financial markets process, and over the last two decades its quantitative study has moved from verbal description to measurement. Entropy and mutual information appear in models of rational inattention \citep{sims2003implications,matejka2015rational,mackowiak2023rational}, in theories of information acquisition and portfolio choice \citep{vannieuwerburgh2010information,veldkamp2011information,kacperczyk2016rational}, in the ordering of signals by their value to investors \citep{cabrales2013entropy,kadan2019estimating}, and in a large empirical literature that measures connectedness, spillovers and the informational content of one market for another \citep{billio2012econometric,diebold2014network}. The Kullback--Leibler divergence \citep{kullback1951information} and mutual information \citep{shannon1948mathematical,cover2006elements} have become the natural units in which such quantities are reported. Whenever the variables are numerous, these units are computed from covariance matrices, typically through the Gaussian formula $I=\tfrac12\log[\det\Sigma_{11}\det\Sigma_{22}/\det\Sigma]$ or through closely related log-determinant expressions. The covariance matrices are themselves estimated, and in modern financial applications the number of variables is rarely small compared with the number of observations.

Random matrix theory describes what happens to such estimates when dimension and sample size are of the same order. Since the work of \citet{laloux1999noise} and \citet{plerou1999universal}, it has been a central tool of econophysics. Those papers showed that the bulk of the eigenvalues of an empirical correlation matrix of stock returns is close to the Marchenko--Pastur law of pure noise \citep{marchenko1967distribution}, that a handful of eigenvalues escape the bulk, and that the largest of them describes a market mode. Subsequent work established the statistical theory of these outliers \citep{johnstone2001distribution,baik2005phase,paul2007asymptotics,benaych2011eigenvalues}, developed estimators that clean the noise from the spectrum \citep{ledoit2004well,ledoit2012nonlinear,bun2016rotational,bun2017cleaning,ledoit2020analytical}, and quantified the risk underestimation of optimised portfolios \citep{pafka2003noisy,elkaroui2010high,bodnar2018estimation}. The monographs of \citet{bai2010spectral}, \citet{anderson2010introduction} and \citet{potters2020first} collect the mathematics. Yet the random matrix theory of \emph{information functionals} of sample covariance matrices, as distinct from their spectra, eigenvectors or portfolio risk, has received far less attention in finance. That gap matters because the information functionals are exactly the quantities that economic theory gives a price to.

This paper develops that theory and applies it. The organising object is the plug-in Gaussian mutual information between two blocks of variables, of dimensions $p$ and $r$, estimated from $n$ observations. I show that its expected excess over the truth, which I call the \emph{spurious information}, is an exact finite-sample function of $(p,r,n)$ alone and is unaffected by the true covariance matrix. It is a telescoping sum of digamma differences, so that its positivity and its strict monotonicity in each block dimension are elementary consequences of the monotonicity of the digamma function. Under independence of the blocks, the plug-in information is minus one half the logarithm of Wilks' statistic. Its exact law is therefore that of a sum of independent log-Beta variables, which gives the exact null variance and an exact test. In the proportional regime $p/n\to c_1$, $r/n\to c_2$, the spurious information grows linearly in $n$ at the rate $\kappa(c_1,c_2)$ displayed in the abstract, while its standard deviation stays of order one. A naive count of parameters would predict the classical value $pr/(2n)$, which underestimates the spurious information by a factor that is already close to two when $c_1=c_2=0.4$.

The same quantity has a second, purely spectral description. The plug-in information equals $-\tfrac12\sum_k\log(1-\hat\rho_k^2)$, where $\hat\rho_k$ are the sample canonical correlations. Under independence their empirical distribution converges to the law of \citet{wachter1980limiting}, the limiting spectral law of the Jacobi or MANOVA ensemble \citep{johnstone2008multivariate,collins2005product}. Matching the two descriptions yields a closed-form evaluation of the logarithmic moment of the Wachter law, and the same argument identifies the variance of the corresponding linear spectral statistic. I also use the finite-rank theory of canonical correlations \citep{bao2019canonical} to describe a phase transition for information. A single dependence direction with population canonical correlation $t$ produces a separated sample canonical correlation only if $t^4>c_1c_2/[(1-c_1)(1-c_2)]$. The corresponding mutual information, $I_c(c_1,c_2)=-\tfrac12\log(1-t_c^2)$, is a frontier below which information exists but cannot be located by spectral methods.

These results translate directly into economics. In a price-taking CARA--Gaussian portfolio problem, I show that the ratio of ex-ante expected utilities with and without a jointly Gaussian signal is exactly $\exp\{-I(v;s)\}$, where $v$ is the vector of payoffs and $s$ the signal. The certainty-equivalent value of information is therefore $I(v;s)/\gamma$, mutual information measured in units of risk tolerance. The identity is in the spirit of \citet{grossman1980impossibility}, \citet{vannieuwerburgh2010information} and \citet{cabrales2013entropy}, but here it holds as an exact equality for arbitrary dimensions. It implies that any econometrician, investor or data vendor who estimates the value of a high-dimensional signal set from data inherits the spurious information as a \emph{spurious value of information}, with the same exact expectation, null law and proportional limit. Because the spurious value grows strictly when uninformative series are appended, a seller can inflate the measured value of a signal set without adding information. A buyer can remove the inflation exactly with a digamma correction. The same exact theory yields the expected Kullback--Leibler cost of acting on plug-in Gaussian beliefs in either direction. These costs are Stein's loss and the entropy loss in information-theoretic form \citep{james1961estimation,dey1985estimation}, and I give their finite-sample values and high-dimensional limits.

Financial correlations are not static, and econophysics has emphasised their regime-like dynamics \citep{munnix2012identifying,schmitt2013non,allez2012eigenvector}. I therefore study the information contained in a continuously observed matrix Ornstein--Uhlenbeck process, whose eigenvalues perform a Dyson Brownian motion with confinement \citep{dyson1962brownian,chan1992wigner,rogers1993interacting}. Its Stieltjes transform obeys a forced complex Burgers equation whose stationary solution is the semicircle. Girsanov's theorem gives the Kullback--Leibler rate between two mean-reversion regimes as $N(N+1)(b-a)^2/(8a)$. The same rate is obtained from the eigenvalue process alone, so the Dyson process is informationally sufficient for the regime and the eigenvector motion carries no information about it. By the Chernoff--Stein lemma, the time needed to tell two regimes apart therefore falls like $N^{-2}$. This sufficiency result connects to variational inference for latent stochastic differential equations \citep{li2020scalable,rice2026sldi}, where the path-space Kullback--Leibler divergence is the regulariser. For rotation-equivariant drift families, inference on the eigenvalue process loses nothing.

The empirical application uses \NDays{} daily value-weighted returns of the \NAssets{} Fama--French industry portfolios from \SampleStart{} to \SampleEnd{} \citep{fama1997industry,french2026data}. The data reproduce the familiar econophysics picture of a Marchenko--Pastur bulk with a dominant market mode and a few sector modes. Inverting the BBP outlier map lets me attach to each sector eigenvector a predicted overlap with its population counterpart, which distinguishes the reliable modes from those that are barely above the noise. Over \NWindows{} rolling windows the market mode absorbs between \AbsMin{} and \AbsMax{} of total variance. A minimum-variance backtest exhibits the noise-dressing penalty predicted by random matrix theory and the benefit of cleaning, which is large when the estimation window is short. The information analysis asks how much the returns of eight carbon-intensive industries reveal about four financial industries, a question motivated by the climate-finance literature on transition risk \citep{battiston2017climate,bolton2021investors,giglio2021climate}. The exact spurious information for this design is \MIBias{} nats with a null standard deviation of \MISd{}. In a time-shifted placebo, the plug-in information averages \MINullPlugMean{} nats, the corrected placebo averages \MINullCorrMean{}, and the corrected test rejects in \MINullRej\% of windows at the nominal 5\% level. The genuine linkage averages \MICorrMean{} nats after correction and is strongly associated with the market mode. Conditioning on the market leaves a smaller linkage of \MIPartialMean{} nats on average. It is significant in every window and uncorrelated with the absorption ratio, and it peaks in \MIPartialMaxDate{}.

The paper draws on, and speaks to, several literatures. Within random matrix theory it builds on the classical theory of Gaussian and Wishart ensembles \citep{mehta2004random,forrester2010log,dumitriu2002matrix,edelman2005random}, on the local laws and anisotropic estimates that make the outlier theory uniform \citep{knowles2017anisotropic}, and on the handbook and textbook treatments of \citet{akemann2011oxford}, \citet{livan2018introduction} and \citet{couillet2022random}. The latter emphasises, as this paper does, that random matrix corrections are first order in modern data analysis. Optimal eigenvalue shrinkage and thresholding in spiked models \citep{gavish2014optimal,donoho2018optimal}, regularisation by banding and thresholding \citep{bickel2008regularized}, factor-based covariance estimation \citep{fan2013large}, and dynamic covariance models with nonlinear shrinkage \citep{engle2002dynamic,engle2019large} form the estimation toolkit against which the information-specific corrections here should be compared. \citet{ledoit2022power} survey that toolkit, and its portfolio applications go back to \citet{ledoit2003improved} and \citet{ledoit2017nonlinear}. In econophysics, the random matrix analysis of financial correlations was extended by \citet{gopikrishnan2001quantifying}, \citet{plerou2002random}, \citet{utsugi2004random} and \citet{burda2005spectral}, and was complemented by hierarchical and network descriptions \citep{mantegna1999hierarchical,tumminello2010correlation}. Textbook accounts include \citet{mantegna2000introduction}, \citet{bouchaud2003theory} and the review of \citet{kwapien2012physical}. The portfolio consequences of estimation noise, from \citet{markowitz1952portfolio} through \citet{elton1973estimating}, \citet{michaud1989markowitz}, \citet{jagannathan2003risk}, \citet{pafka2002noisy} and \citet{kondor2007noise}, motivate the minimum-variance exercise. In financial econometrics, the determination of the number of factors \citep{bai2002determining,onatski2009testing,onatski2010determining,ahn2013eigenvalue} is itself a random matrix problem rooted in the arbitrage pricing theory \citep{ross1976arbitrage,chamberlain1983arbitrage}. So is the estimation of latent factors from large panels and high-frequency data \citep{lettau2020factors,lettau2020estimating,pelger2020understanding,aitsahalia2017using}. Random matrix asymptotics also explain the surprising behaviour of overparameterised regression \citep{hastie2022surprises}. In information theory, random matrix methods have long been used to compute the capacity of multi-antenna channels, which is a log-determinant functional of exactly the kind studied here \citep{telatar1999capacity,tulino2004random}. The difficulty of estimating mutual information in high dimensions is well documented for nonparametric and neural estimators \citep{paninski2003estimation,kraskov2004estimating,belghazi2018mutual,mcallester2020formal}. The Gaussian case studied here is the one in which that difficulty can be computed exactly. In economics, the comparison of information structures goes back to \citet{blackwell1953equivalent}. The role of information in markets is surveyed by \citet{stiglitz2000contributions}, and its incorporation into prices is the subject of the efficiency and microstructure literatures \citep{fama1970efficient,kyle1985continuous}.

The paper makes four contributions. The first is a set of exact and asymptotic laws for spurious information, Kullback--Leibler costs and their fluctuations, derived from the Bartlett decomposition, the Jacobi ensemble and the Marchenko--Pastur and Wachter laws, including a new closed form for the logarithmic moment of the Wachter law. The second is the economic translation of these laws through an exact value-of-information identity, together with the monotonicity, correction and frontier results it implies. The third is an information-theoretic analysis of Dyson dynamics that establishes the sufficiency of the eigenvalue process and the quadratic scaling of regime detectability. The fourth is an econophysics application that combines these tools on more than three decades of data. Every formula is checked against simulation, and the combinatorial and algebraic steps that link the analytic arguments are verified in the Lean~4 theorem prover \citep{demoura2021lean}.

The paper proceeds as follows. The next section reviews the random matrix machinery used throughout, namely resolvents, Stieltjes and $S$-transforms, the Marchenko--Pastur and Wachter laws, spiked models and Dyson Brownian motion. I then derive the information laws and their economic interpretation, followed by the dynamic results. A numerical section verifies the theory, and the application section takes it to the data. The appendices contain proofs, a description of the Lean formalization and additional tables.
